% ============================================================
\section{Tooling \& Testing (Lab 12)}
% ============================================================
\subsection{Unit-Tests mit GoogleTest}
\emph{Unit-Tests} pr\"ufen einzelne Bausteine (z.\,B. eine Klasse) automatisiert. GoogleTest (gtest) ist ein verbreitetes Framework.
\begin{syntaxbox}
\code{TEST(TestSuiteName, TestName) \{ ... Assertions ... \}}
\end{syntaxbox}
\begin{lstlisting}
#include <gtest/gtest.h>
#include "calculator.hpp"

TEST(CalculatorTest, Addition) {
    Calculator calc;
    EXPECT_EQ(calc.add(2, 3), 5);
    EXPECT_EQ(calc.add(-1, 1), 0);
}

TEST(CalculatorTest, DivisionByZero) {
    Calculator calc;
    EXPECT_THROW(calc.divide(10, 0), std::invalid_argument);   // erwartet Exception
}
\end{lstlisting}
\begin{center}\small
\begin{tabular}{@{}ll@{}}
\toprule
\textbf{Assertion} & \textbf{pr\"uft} \\
\midrule
\code{EXPECT\_EQ(a, b)} & \code{a == b} (Test l\"auft bei Fehler weiter) \\
\code{EXPECT\_TRUE(x)} / \code{FALSE} & Wahrheitswert \\
\code{EXPECT\_THROW(expr, Typ)} & dass \code{expr} eine Exception vom \code{Typ} wirft \\
\code{ASSERT\_*} & wie \code{EXPECT\_*}, aber bricht den Test bei Fehler ab \\
\bottomrule
\end{tabular}
\end{center}
\begin{tipp}
\code{EXPECT\_*} meldet einen Fehler, l\"asst den Test aber weiterlaufen (mehrere Pr\"ufungen pro Test). \code{ASSERT\_*} bricht sofort ab -- sinnvoll, wenn die folgenden Zeilen sonst abst\"urzen w\"urden (z.\,B. nach einem Null-Check).
\end{tipp}

\subsection{Build mit CMake}
\emph{CMake} beschreibt den Build plattformunabh\"angig in \code{CMakeLists.txt}. \code{FetchContent} l\"adt Abh\"angigkeiten (hier gtest) automatisch.
\begin{lstlisting}[language={},morekeywords={}]
cmake_minimum_required(VERSION 3.20)
project(gtest_lab LANGUAGES CXX)
set(CMAKE_CXX_STANDARD 17)
set(CMAKE_CXX_STANDARD_REQUIRED ON)

enable_testing()
add_library(calculator src/calculator.cpp)
target_include_directories(calculator PUBLIC ${CMAKE_CURRENT_SOURCE_DIR}/include)

include(FetchContent)
FetchContent_Declare(googletest
    URL https://github.com/google/googletest/archive/refs/tags/v1.17.0.zip)
FetchContent_MakeAvailable(googletest)

add_executable(calculator_tests tests/calculator_tests.cpp)
target_link_libraries(calculator_tests PRIVATE calculator GTest::gtest_main)

include(GoogleTest)
gtest_discover_tests(calculator_tests)
\end{lstlisting}
Bauen \& Testen typischerweise:
\begin{lstlisting}[language={}]
cmake -S . -B build      # Build-Verzeichnis konfigurieren
cmake --build build      # kompilieren
ctest --test-dir build   # Tests ausfuehren
\end{lstlisting}
\begin{merke}
Grundbausteine einer \code{CMakeLists.txt}: \code{add\_library} / \code{add\_executable} (Ziele), \code{target\_link\_libraries} (Abh\"angigkeiten/Bibliotheken), \code{target\_include\_directories} (Header-Pfade). \code{gtest\_discover\_tests} registriert die Tests automatisch bei CTest.
\end{merke}
\clearpage
